\section{Objectifs}
\begin{itemize}
\item Modéliser une évolution à taux constant.
\item Utiliser une suite géométrique pour une évolution discrète.
\item Comprendre l'extension vers la fonction exponentielle $x\mapsto a^x$ pour $a>0$ et $x\ge0$.
\item Calculer un taux d'évolution moyen et résoudre un problème de seuil.
\end{itemize}

\section{1. Évolution exponentielle}
Une évolution exponentielle correspond à un \textbf{coefficient multiplicateur constant}.

Une hausse de 5 \% correspond au coefficient
\[
1+\frac{5}{100}=1{,}05.
\]

Une baisse de 8 \% correspond au coefficient
\[
1-\frac{8}{100}=0{,}92.
\]

\section{2. Suite géométrique}
Une suite géométrique à termes strictement positifs et de raison $q>0$ vérifie
\[
u(n+1)=q\,u(n).
\]

Ainsi,
\[
u(n)=u(0)q^n.
\]

Si $q>1$, la suite est croissante. Si $0<q<1$, elle est décroissante.

\section{3. Fonction exponentielle}
La fonction
\[
x\mapsto a^x\qquad(a>0,\ x\ge0)
\]
prolonge le modèle géométrique aux valeurs réelles positives de $x$.

On utilise notamment
\[
a^{x+y}=a^x a^y.
\]

Les valeurs comme $a^{1/n}$ permettent d'interpréter une évolution sur une fraction de période.

\section{4. Taux d'évolution moyen}
Si un coefficient global $C$ résulte de $n$ évolutions identiques, le coefficient moyen $q$ vérifie
\[
q^n=C.
\]

Le taux moyen vaut alors $q-1$.

Exemple : une quantité est multipliée par 1,21 en deux périodes :
\[
q=\sqrt{1{,}21}=1{,}10.
\]
Le taux moyen est donc 10 \% par période.

\section{5. Seuil et ordre de grandeur}
Pour déterminer quand une quantité dépasse un seuil, on peut calculer des termes successifs, utiliser une représentation graphique ou un outil numérique.

Une croissance exponentielle peut devenir très rapide : l'ordre de grandeur du résultat doit toujours être contrôlé.

\section{6. Applications}
\begin{itemize}
\item intérêts composés ;
\item décroissance radioactive ;
\item élimination d'une substance ;
\item croissance de populations selon un modèle simplifié.
\end{itemize}

\section{Automatismes à entretenir}
\begin{itemize}
\item Passer d'un taux au coefficient multiplicateur.
\item Calculer des évolutions successives.
\item Utiliser les puissances.
\item Lire une croissance ou une décroissance sur un graphique.
\end{itemize}

\section{Erreurs fréquentes}
\begin{itemize}
\item Confondre ajout constant et pourcentage constant.
\item Additionner des pourcentages successifs au lieu de multiplier les coefficients.
\item Utiliser une raison négative alors que le programme traite ici des suites géométriques à termes strictement positifs.
\end{itemize}

\section{À retenir}
Évolution exponentielle = coefficient multiplicateur constant. En discret : suite géométrique. Sur les réels positifs : fonction exponentielle.
